\noindent
\graphicspath{{./21-08/Images/}} 
%%%%%%%%%%%%%%%%
% Page de gauche
%%%%%%%%%%%%%%%%
\begin{minipage}[t][\textheight][t]{0.33\textwidth}
    \colorbox{StrongBlue}{
        \parbox[c][\textheight][c]{\textwidth}{ 
            \color{OffWhite} 
            \begin{center}
                \begin{minipage}[t][0.6\textheight][c]{0.8\textwidth}
                    {\huge \textbf{Jeudi 21 Août}}\\

                    \normalsize
                    \begin{center}
                        \begin{tabular}{@{} l @{\hspace{0.5cm}} l @{}}
                            Abri de départ : &\hfill \makecell{Saint Gildas ou,\\ Beg Run er Vilin } \\ 
                            Abri d'arrivée : &\hfill Port Haliguen \\ 
                            Distance : &\hfill 13 Milles \\ 
                            Départ : &\hfill 09 h 30 \\ 
                            Arrivée prévue : &\hfill 17 h 00 \\ 
                            Lever soleil 
\includegraphics[width=8px,height=8px]{../../Pictos/soleil_OffWhite.png} : &\hfill 07 h 16 \\ 
                            Coucher soleil 
\includegraphics[width=8px,height=8px]{../../Pictos/lune_OffWhite.png} : &\hfill 21 h 11 \\ 
                            \\
                            Abris replis : &  \\
                                    & • Saint Gildas \\
                                    & • Crouesty \\
                                    & • La Trinité \\
                                    & • Port Haliguen \\                    
                                
                        \end{tabular}
                    \end{center}
                    
                \end{minipage}%

                \begin{minipage}[t][0.4\textheight][t]{\textwidth}

                    \begin{center}
    {\large \textit{Marées \textbf{Port Haliguen}}}\\ \vspace{0.5em}
    
    \renewcommand{\arraystretch}{1.40} % Adjust the row height by a factor of 1.4
    
    \begin{tabular}{ |C{5em}|c|C{4em}|C{4em}| } \hline
        & Coef & Heure & Hauteur\\ \hline
        BM & \multirow{1}{1em}{67}  & 10 h 23 & 1.37 m\\ \hline 
        PM & \multirow{2}{1em}{73}  & 16 h 52 & 4.85 m\\ \cline{1-1} \cline{3-4}
        BM &   				        & 04 h 36 & 4.61 m\\ \hline
    \end{tabular}
\end{center}

                    \begin{center}
    {\large \textit{Marées port de \textbf{saint Gildas}}}\\ \vspace{0.5em}    
    
    \renewcommand{\arraystretch}{1.40} % Adjust the row height by a factor of 1.5
    
   \begin{tabular}{ |C{5em}|c|C{4em}|C{4em}| } \hline
        & Coef & Heure & Hauteur\\ \hline
        PM & \multirow{2}{1em}{54}  & 04 h 33 & 4.91 m\\ \cline{1-1} \cline{3-4}
        BM &                        & 10 h 14 & 1.54 m\\ \hline
        PM & \multirow{1}{1em}{60}  & 16 h 44 & 5.23 m\\ \hline
    \end{tabular}
\end{center}
                \end{minipage}%
            \end{center}       
        }      
    }
\end{minipage}%
\begin{minipage}[c][\textheight][t]{0.10\textwidth}
\hfill
\end{minipage}%
\begin{minipage}[c][\textheight][t]{0.50\textwidth}
    \vspace{4cm}
    \color{StrongBlue}    
    {\huge \textbf{Description de la navigation}}\\
            \begin{itemize}
                \item La navigation se fera en équipe de nuit, briefing de navigation à 8h30 sur un bateau de rendez-vous. Les annexes sont rammenées à leur bateau d'origine et dégonflées à l'exception de celle du bateau de rendez-vous qui servira à rammener les PE à leurs bords.  
                \item Départ 9h30. Cap plein Nord afin de rejoindre la cardinale Sud \textit{Buissons de Méhaban}. 
                \item Organisation d'une régate côtière en mobilisant les cardinales et dangers isolés présents pour matérialiser le parcours en fonction du vent~: Sud \textit{Buissons de Méhaban}, Ouest \textit{Roche de la révision}, cardinale \textit{Quiberon Nord}, dangers isolés \textit{Bugalet} et \textit{Men er Roué}.
                \item Repas en navigation, avant ou après la régate en fonction de la durée du trajet.  
                \item Une fois la régate terminée, retour au port. Douche, repas puis veillée sur la plage. 
                \item Un réassort en produits frais est possible 
            \end{itemize}
        {
        \huge \textbf{Contraintes de la navigation}}\\
            \begin{itemize}
                \item Au départ, rappel des étappes nécessaires pour lever un mouillage. 
                \item On veillera à ne pas aller dans \textit{Les Buissons de Méhaban}.
                \item La marée sera montante ou haute lors de l'arrivée au port, la cardinale Sud \textit{Port Haliguen} n'est donc pas un obstacle. On veillera néanmoins à respecter la \textit{S.E. Olibarte} et la cardinale Nord adjacente. 
            \end{itemize}
\end{minipage}%
%%%%%%%%%%%%%%%%
% Page de droite
%%%%%%%%%%%%%%%%
\newgeometry{top=2cm, bottom=2cm, left=2cm, right=2cm}
\begin{minipage}[t][0.9\textheight][t]{0.60\textwidth}
    \color{StrongBlue}    
    \noindent
    \begin{minipage}[t][0.3\textheight][t]{\textwidth}
        \begin{minipage}[t][0.3\textheight][t]{0.65\textwidth}
            \vspace{-8.5em}
            {\huge \textbf{Détails de la navigation}}\\
            Le QR-code situé à droite est cliquable, il permet d'explorer le planning de navigation sur le site datashom. 
            La carte ci-dessous visualise le trajet prévu ainsi que les abris de replis et dangers identifiés. 
        \end{minipage}%
        \hfill
        \begin{minipage}[t][0.3\textheight][t]{0.25\textwidth}       
        \href{https://data.shom.fr/datacontext/fr/1747929590690-267863}{ 
            \includegraphics[width=0.8\textwidth]{qr-code.png}
        }
        \end{minipage}%
    \end{minipage}
    \begin{minipage}[t][0.7\textheight][t]{\textwidth}
        % \includegraphics[width=0.97\textwidth]{carte.png}
        \includegraphics[height=0.7\textheight]{carte.png}
    \end{minipage}
\end{minipage}
\begin{minipage}[t][0.9\textheight][t]{0.40\textwidth}
    \color{StrongBlue} 
    \centering 
    \includegraphics[height=\textheight, width=\textwidth, keepaspectratio]{illustration_verticale.jpg}
\end{minipage}%
\restoregeometry



